\section{The critical singularity and both limiting directions}\label{sec:critical}

The critical formula~\eqref{eq:products-main} yields
\begin{equation}\label{eq:log-derivative}
 \frac{D_1'(u)}{D_1(u)}=-\frac1{2(1-u)}-\frac{\log(1-u)}{2u},
 \qquad |u|<1.
\end{equation}
Consequently
\begin{equation}\label{eq:fractional-order}
 \lim_{u\uparrow1}(u-1)\frac{D_1'(u)}{D_1(u)}=\frac12.
\end{equation}
A nonzero meromorphic function near one has the form $(u-1)^k h(u)$, with
$k\in\mathbb Z$ and $h(1)\ne0$ holomorphic, and the same limit would be the
integer $k$. This proves the non-meromorphic assertion. It excludes an entire
ordinary trace-class Fredholm determinant with this germ. It does not give a
natural boundary on the whole unit circle: the logarithm and dilogarithm admit
compatible local continuations near its other points.

For fixed $n$, binomial expansion gives
\begin{align}
 \frac1{1-(1-h)^n}&=\frac1{nh}+\frac{n-1}{2n}+O_n(h),\label{eq:minus-exp}\\
 \frac1{(1+h)^n-1}&=\frac1{nh}-\frac{n-1}{2n}+O_n(h).\label{eq:plus-exp}
\end{align}
The first formula with $a=1-\varepsilon$, $q=1-2\varepsilon$ shows
$\tau_n(a)\to(n-1)/(2n)$ as $a\uparrow1$. A uniform bound is still needed to
pass through the full logarithmic series. For $1/2<a<1$, the relation
$q\le a^2$ follows from $a^2-q=(a-1)^2$. Convexity of $x^n$ on the positive
line, with $a=(1+q)/2$, gives $1-2a^n+q^n\ge0$. Therefore
\begin{equation}\label{eq:uniform-minus}
 0\le\tau_n(a)
 \le\frac2{1-a^{2n}}-\frac1{1-a^n}
 =\frac1{1+a^n}\le1.
\end{equation}

For $a>1$, the two finite fixed points $\pm\sqrt{b/(a-1)}$ both have
multiplier $q$. Formula~\eqref{eq:primitive} shows that their full primitive
factor is precisely $F_a$ from~\eqref{eq:reduced-definition}. Removing them
means subtracting their contributions at \emph{every} repetition:
\begin{equation}\label{eq:reduced-tau}
 \tau_n^{\red}(a)=\frac1{a^n-1}-\frac2{q^n-1},\qquad
 D_a^{\red}(u)=\exp\left(-\sum_{n\ge1}\frac{\tau_n^{\red}(a)}n u^n\right).
\end{equation}
The factor $F_a$ is nonzero on $|u|<1$, so the quotient is unambiguous.
Using~\eqref{eq:plus-exp} with $a=1+\varepsilon$, $q=1+2\varepsilon$ proves
$\tau_n^{\red}(a)\to(n-1)/(2n)$. Convexity gives
$q^n-2a^n+1\ge0$, and $q\le a^2$ yields the uniform bound
\begin{equation}\label{eq:uniform-plus}
 0\le\tau_n^{\red}(a)
 \le\frac1{a^n-1}-\frac2{a^{2n}-1}
 =\frac1{1+a^n}\le\frac12.
\end{equation}

Fix $\rho<1$. The summands of both normalized logarithmic series are bounded
on $|u|\le\rho$ by $\rho^n/n$, a summable sequence independent of the
approaching parameter. Fixed-coefficient convergence and a finite-head,
uniform-tail argument give uniform convergence of the logarithms.
Exponentiation preserves it on this compact disk, proving both limits in
Theorem~\ref{thm:critical}.

The uncompensated limit is different. For real $0<u<1$,
\begin{equation}\label{eq:unreduced}
 \log D_a(u)\le-\tau_1(a)u=-\frac u{a-1}\longrightarrow-\infty
 \quad(a\downarrow1),
\end{equation}
whereas $D_a(0)=1$. It cannot converge locally uniformly to $D_1$, which is
nonzero on the open disk. The required compensation is fixed by the two
real orbits colliding with infinity, not by matching the desired limit after
an arbitrary scalar normalization.
